\subsection{Systematics of Moduli Stabilisation in Calabi-Yau Flux Compactifications --- arXiv:hep-th/0502058}

\textbf{Authors:} Vijay Balasubramanian, Per Berglund, Joseph P. Conlon, Fernando Quevedo

\paragraph{主な主張}
Type IIBのCalabi--Yauフラックスコンパクト化で、\(\alpha'\)補正と非摂動効果の競合により、指数関数的に大きな体積を持つ非超対称AdS極小が生じ得ることを示した。 \cite{Balasubramanian:2005zx}

\paragraph{新規性}
大体積でのポテンシャルの系統的な次数評価から、摂動的補正を不可欠な安定化要素とする真空を見いだし、具体的なCalabi--Yauで実証した。

\paragraph{位置付け}
Large Volume Scenario（LVS）の原論文であり、大体積を用いたモジュライ固定とスケール階層の構成の基礎である。

\paragraph{コメント（読書メモ）}
Large Volume Scenario（LVS）の原論文として整理する。
